\begin{landscape}
\begin{figure}[tbp]
    \centering
    \includegraphics[width=\linewidth,height=0.73\textheight,keepaspectratio]{jaxka-gap-cutoffs/fig.pdf}
    \caption[Residue-deletion calibration of gap masks]{Residue-deletion tests used to calibrate gap masks. Panels separate terminal and internal deletions by gap length and predictor. Blue curves show changes in binding-induced shifts; orange curves show the larger absolute change in bound- or separated-state $\mathrm{p}K_a$. Grey density shows absolute changes in binding-induced shifts, with exact zeros in a separate strip below the logarithmic axis and nonzero changes below $10^{-3}$ plotted at $10^{-3}$. Curves show medians and 95th percentiles in 5~\AA{} distance bands. Short horizontal marks give upper 95\% bootstrap bounds on the 95th percentile from 1,000 complex-level resamples. Bounds require at least ten complexes and twenty observations. The horizontal reference is the 0.1 $\mathrm{p}K_a$ tolerance. Short-gap distances use observed flanking backbone atoms, while long-gap distances use flanking C$\alpha$ atoms. PypKa tests cover 23 complexes and pKAI tests 515; pKAI reports $\mathrm{p}K_a$ to 0.01. Resampling whole complexes retains the dependence between repeated site observations. Production masks combine these calibrations with the anchor and clearance rules in \cref{tab:jaxka-gap-masks}.}
    \label{fig:jaxka-gap-cutoffs}
\end{figure}
\end{landscape}
